% GENERATED FILE. Edit canonical lessons, then run npm run generate:books.
% canonical-source-hash: fnv1a64:4cacc8aa4dbfcb69
% canonical-lessons: LA-C14-annos-natus

\chapter{Asking Someone's Age}
\label{ch:la-asking-age}

This chapter is generated from the canonical micro-lessons used by Language Ladder.
Each section stays independently resumable and preserves the authored prerequisite order.

\section[vīgintī annōs nātus sum]{vīgintī annōs nātus sum — a THIRD way to say your age}

\label{lesson:LA-C14-annos-natus}



\begin{quote}
\textbf{Warm-up.} {[}PAUSE 2s{]} You already know that Spanish and French say \textquotedblleft{}I HAVE my years,\textquotedblright{} and German/English say \textquotedblleft{}I AM my years.\textquotedblright{} Here's the honest surprise: Latin itself did \textbf{neither}.
\end{quote}

\subsection*{You'll want to know: vīgintī annōs nātus sum — \textquotedblleft{}twenty years born\textquotedblright{}}
\textbf{vīgintī annōs nātus sum} = \textquotedblleft{}\textbf{I am twenty years old}\textquotedblright{} — literally \textquotedblleft{}\textbf{I am, having been born, twenty years}\textquotedblright{} (a man speaking; a woman says \textbf{nāta}). Taken apart:

\begin{itemize}
\raggedright
  \item \textbf{nātus/nāta} — \textquotedblleft{}\textbf{born},\textquotedblright{} the perfect passive participle of \textbf{nāscor}, \textquotedblleft{}to be born\textquotedblright{} (the same root behind English \textbf{nature, native, nation, innate}).
  \item \textbf{annōs} — \textquotedblleft{}\textbf{years},\textquotedblright{} in the \textbf{accusative case}, marking \textbf{duration} — \textquotedblleft{}for twenty years {[}of life so far{]}.\textquotedblright{} This is the same accusative-of-duration you'd use for \textquotedblleft{}I studied \textbf{for three years}\textquotedblright{} (\emph{trēs annōs}).
  \item \textbf{sum} — \textquotedblleft{}I am.\textquotedblright{}
\end{itemize}

So the whole phrase reads: \textquotedblleft{}I am — having been born {[}and living{]} twenty years {[}so far{]}.\textquotedblright{} Cicero uses exactly this pattern in \emph{Cato Maior de Senectute} (§14), in Cato's own voice: \textbf{quīnque et sexāgintā annōs nātus}, \textquotedblleft{}having been born sixty-\textbf{five} years,\textquotedblright{} i.e., \textquotedblleft{}at sixty-five years old\textquotedblright{} — the same section also gives Ennius's age as \textbf{annōs septuāgintā nātus}, \textquotedblleft{}seventy.\textquotedblright{}

\begin{culture}
You've already met the clean Romance/Germanic split for this concept — Romance languages \textbf{have} their years (\emph{tengo}, \emph{j'ai}, \emph{ho}), Germanic languages \textbf{are} their years (\emph{ich bin}, \emph{I am}). Here's the twist: \textbf{Latin itself used neither construction.} Classical Latin's \emph{nātus} + accusative-of-duration pattern is a \textbf{third}, distinct grammatical shape — not \textquotedblleft{}having\textquotedblright{} years, not simply \textquotedblleft{}being\textquotedblright{} a number, but \textquotedblleft{}being \textbf{born-for-that-long}.\textquotedblright{} The Romance \textquotedblleft{}have your years\textquotedblright{} pattern (\emph{habēre} + years) that Spanish, French, and Italian all inherited is a \textbf{Vulgar Latin innovation} that arose \textbf{after} Classical Latin, replacing the older \emph{nātus} construction — so Spanish's \emph{tengo veinte años} is NOT a direct continuation of how Cicero would have said it.
\end{culture}

\subsection*{Guided Practice}
{[}PAUSE 1s{]}

\begin{itemize}
\raggedright
  \item {[}YOU SAY: \textquotedblleft{}vīgintī annōs nātus sum\textquotedblright{} — I am twenty years old, literally \textquotedblleft{}born twenty years\textquotedblright{}{]}
  \item {[}YOU SAY: \textquotedblleft{}nātus\textquotedblright{} — born, same root as English \textquotedblleft{}native, nature, nation\textquotedblright{}{]}
  \item {[}YOU SAY: the honest surprise — this is a THIRD pattern, neither \textquotedblleft{}have\textquotedblright{} nor simple \textquotedblleft{}be\textquotedblright{}{]}
\end{itemize}

\begin{tcolorbox}[breakable,colback=teal!4,colframe=teal!35!black,title={Wrap-up recall}]
{[}PAUSE 3s{]} What does \textbf{vīgintī annōs nātus sum} literally mean? (\textquotedblleft{}I am, having been \textbf{born}, twenty years\textquotedblright{} — using the perfect participle \emph{nātus} + the accusative of duration.) Does this match Romance's \textquotedblleft{}have your years\textquotedblright{} pattern? (\textbf{No} — that's a \textbf{later Vulgar Latin innovation}, not what Classical Latin did.) What English words share the root of \emph{nātus}? (\textbf{Native, nature, nation, innate}.)
\end{tcolorbox}
